% !TEX TS-program = pdflatex
% !TEX encoding = UTF-8 Unicode
%
% Resumen (2 páginas) del Capítulo 1 — "Introducción"
% Libro: Inteligencia Computacional: Fundamentos y Aplicaciones
% Autor: Alexandre G. Evsukoff (E-papers, Rio de Janeiro, 2020)
%
% Idioma: Español (español latinoamericano/chileno académico).
% Generado para el curso de Inteligencia Computacional — USACH (Prof. Max Chacón).

\documentclass[11pt,a4paper]{article}

\usepackage[utf8]{inputenc}
\usepackage[T1]{fontenc}
\usepackage[spanish]{babel}
\usepackage[margin=2.2cm]{geometry}
\usepackage{microtype}
\usepackage{parskip}
\usepackage{enumitem}
\usepackage{hyperref}
\hypersetup{colorlinks=true, linkcolor=black, urlcolor=blue!60!black}

\setlist[itemize]{leftmargin=1.2em, itemsep=1pt, topsep=2pt}

\title{\vspace{-1.5cm}Resumen --- Capítulo 1: Introducción\\
\large Inteligencia Computacional: Fundamentos y Aplicaciones \\
\normalsize Alexandre G.\ Evsukoff (E-papers, 2020)}
\author{}
\date{}

\begin{document}
\maketitle
\vspace{-1.2cm}

\section*{1. Contexto histórico y definición}

La \textbf{inteligencia computacional} (IC) se define como un conjunto de
paradigmas computacionales destinados al desarrollo y aplicación de sistemas
inteligentes para resolver problemas complejos del mundo real. Su historia está
vinculada a la de la inteligencia artificial (IA): desde las reflexiones filosóficas
sobre la mente y el lenguaje hasta el \emph{test de Turing} (1950). En términos
metodológicos se distinguen dos corrientes principales: la corriente
\emph{simbolista} (sistemas expertos, lógica) y la corriente
\emph{conexionista} (redes neuronales, iniciadas con el Perceptron de Rosenblatt,
1958). La incapacidad del Perceptron para resolver problemas no linealmente
separables llevó a una fase de estancamiento hasta que, en 1986, la técnica de
\emph{retropropagación del error} (Rumelhart, Hinton y Williams) revitalizó las
redes neuronales multicapa. Paralelamente, Quinlan introdujo el algoritmo \textbf{ID3}
para árboles de decisión, consolidando el aprendizaje automático simbólico.

Durante las décadas de 1980 y 1990 surgieron otros paradigmas (lógica difusa,
algoritmos evolutivos, máquinas de soporte vectorial). La integración de estas
técnicas recibió nombres como \emph{soft computing} o \textbf{inteligencia
computacional}. Desde 2010 el término \emph{big data} ganó relevancia, y a partir de
2012 las \textbf{redes neuronales profundas} (\emph{deep learning}) ejecutadas en GPUs
lograron rendimientos comparables o superiores al humano en tareas como visión
computacional y procesamiento de lenguaje natural.

\section*{2. Ciencia de datos y el proceso KDD}

La \textbf{ciencia de datos} se sitúa en la intersección entre matemática,
estadística, informática, ingeniería y conocimiento de dominio. A diferencia del
método científico clásico, en ciencia de datos los modelos se construyen a partir
de los datos observados en vez de partir de hipótesis teóricas predefinidas. El autor
subraya que todo modelo es una representación parcial del fenómeno real y que su
validez depende de la calidad y representatividad de la muestra.

El proceso de \textbf{descubrimiento de conocimiento en bases de datos (KDD)} —también
llamado minería de datos— es cíclico y comprende las etapas:
\begin{enumerate}[leftmargin=1.4em, itemsep=1pt, topsep=2pt]
  \item \textbf{Selección} --- elegir subconjuntos y variables relevantes (conjunto de entrenamiento).
  \item \textbf{Limpieza} --- tratar valores ausentes y \emph{outliers}.
  \item \textbf{Preprocesamiento} --- normalizaciones y transformaciones (p.ej., ACP/PCA).
  \item \textbf{Modelado} --- ajuste de modelos y parámetros sobre los datos de entrenamiento.
  \item \textbf{Evaluación} --- validación estadística de los resultados.
\end{enumerate}
Participación del experto en dominio es esencial en todas las fases.

\section*{3. Tipos de modelos}

Un modelo es una representación simplificada del proceso real. Según su origen,
se distinguen \textbf{modelos de conocimiento} (basados en leyes o relaciones físicas)
y \textbf{modelos de comportamiento} (input-output, ajustados a partir de datos),
siendo la mayoría de las técnicas de IC parte de esta segunda categoría.

Según su propósito, los problemas se dividen en \textbf{preditivos} y
\textbf{descriptivos}: los primeros buscan predecir variables objetivo (aprendizaje
supervisado), y los segundos describir estructuras o patrones sin variable objetivo
(aprendizaje no supervisado). Existe además el \emph{aprendizaje por reforzamiento},
donde un agente aprende mediante interacción y retroalimentación del entorno.

\textbf{Regresión.} Predicción de una variable numérica a partir de atributos de
entrada. La regresión lineal es el punto de partida y se extiende a modelos para
series temporales y sistemas dinámicos. La evaluación suele basarse en medidas de
error de predicción.

\textbf{Clasificación.} Asignación de una etiqueta discreta a partir de atributos;
se busca una función discriminante que separe las clases en el espacio de entrada.
Las clases pueden ser linealmente separables o no, lo que determina la complejidad
del modelo requerido. Aplicaciones típicas: diagnóstico médico, evaluación crediticia,
bioinformática y clasificación de documentos e imágenes.

\textbf{Clustering (agrupamiento).} Identificación de grupos de objetos similares sin
una variable objetivo. Evaluar la calidad de un clustering es complejo; existen
medidas internas y externas, y en muchos métodos el número de grupos es un parámetro
de entrada. Usos frecuentes: segmentación de clientes, análisis genómico y visión por
computador.

Se mencionan también otros tópicos: \emph{reglas de asociación}, aprendizaje de
representaciones, clasificación semisupervisada y métodos de \emph{ensemble}, muy
útiles en competencias de ciencia de datos.

\section*{4. Recursos complementarios}

\textbf{Bibliografía.} Para fundamentos estadísticos y de aprendizaje: Hastie,
Tibshirani y Friedman (\emph{The Elements of Statistical Learning}); Bishop (\emph{Pattern
Recognition and Machine Learning}); Duda \& Hart; Theodoridis \& Koutroumbas. Para redes
neuronales: Haykin; para lógica difusa: Pedrycz \& Gomide; para SVM: Schölkopf \& Smola
y Vapnik. En minería de datos: Han \& Kamber, Tan et al., Aggarwal y Witten \& Frank.
En \emph{deep learning}: Goodfellow, Bengio \& Courville.

\textbf{Recursos en línea.} Portales y repositorios: KDnuggets, UCI Machine Learning
Repository, KEEL y Kaggle. Cursos abiertos (Coursera, Stanford), preprints en arXiv y
código en GitHub. Conferencias relevantes: KDD, ICDM, NeurIPS, ICML, entre otras.

\textbf{Software.} Comerciales: SAS, IBM SPSS Modeler, Statistica; en entorno libre:
Weka, RapidMiner, KNIME, Orange. Lenguajes y bibliotecas: Python (pandas,
scikit-learn, NumPy, SciPy), R y Matlab. KDnuggets muestra el liderazgo de R y Python
entre practicantes.

\section*{5. Síntesis y relevancia para el curso}

El Capítulo 1 establece el vocabulario y el marco conceptual requerido para el resto
del libro: diferencia entre dato/información/conocimiento, el ciclo KDD, y la
taxonomía de modelos (preditivos vs. descriptivos; supervisado vs. no supervisado).
Se destaca la idea crítica de que todo modelo depende de la muestra y de las
suposiciones del proceso de muestreo.

Este planteamiento coincide con la \textbf{Unidad 1 (Introducción)} del programa de
Inteligencia Computacional del Magíster USACH (Prof. Max Chacón) y prepara para los
temas evaluados en la PEP1 (ACP, reglas de asociación, clustering) y las unidades
posteriores (clasificación bayesiana, árboles de decisión, paradigma conexionista).

\vspace{0.4em}
\hrule
\vspace{0.3em}
{\footnotesize Fuente: Evsukoff, A.\ G.\ \emph{Inteligencia Computacional: Fundamentos y
Aplicaciones}. Rio de Janeiro: E-papers, 2020 (Presentación y Capítulo 1, pp. 8--29).
Resumen elaborado para fines académicos en el contexto del curso de Inteligencia
Computacional del Magíster en Ingeniería, Universidad de Santiago de Chile (USACH).}

\end{document}
